% Options for packages loaded elsewhere
\PassOptionsToPackage{unicode}{hyperref}
\PassOptionsToPackage{hyphens}{url}
\documentclass[
]{IEEEtran}
\usepackage{xcolor}
\usepackage{amsmath,amssymb}
\setcounter{secnumdepth}{-\maxdimen} % remove section numbering
\usepackage{iftex}
\ifPDFTeX
  \usepackage[T1]{fontenc}
  \usepackage[utf8]{inputenc}
  \usepackage{textcomp} % provide euro and other symbols
\else % if luatex or xetex
  \usepackage{unicode-math} % this also loads fontspec
  \defaultfontfeatures{Scale=MatchLowercase}
  \defaultfontfeatures[\rmfamily]{Ligatures=TeX,Scale=1}
\fi
\usepackage{lmodern}
\ifPDFTeX\else
  % xetex/luatex font selection
\fi
% Use upquote if available, for straight quotes in verbatim environments
\IfFileExists{upquote.sty}{\usepackage{upquote}}{}
\IfFileExists{microtype.sty}{% use microtype if available
  \usepackage[]{microtype}
  \UseMicrotypeSet[protrusion]{basicmath} % disable protrusion for tt fonts
}{}
\makeatletter
\@ifundefined{KOMAClassName}{% if non-KOMA class
  \IfFileExists{parskip.sty}{%
    \usepackage{parskip}
  }{% else
    \setlength{\parindent}{0pt}
    \setlength{\parskip}{6pt plus 2pt minus 1pt}}
}{% if KOMA class
  \KOMAoptions{parskip=half}}
\makeatother
\setlength{\emergencystretch}{3em} % prevent overfull lines
\providecommand{\tightlist}{%
  \setlength{\itemsep}{0pt}\setlength{\parskip}{0pt}}
\usepackage{bookmark}
\IfFileExists{xurl.sty}{\usepackage{xurl}}{} % add URL line breaks if available
\urlstyle{same}
% fallback for those not using the hyperref driver hyperxmp:
\makeatletter
\@ifundefined{xmpquote}{\newcommand{\xmpquote}[1]{#1}}{}
\makeatother
\hypersetup{
  hidelinks,
  pdfcreator={LaTeX via pandoc}}

\author{}
\date{}

\begin{document}

\section{Design of a Neural Collaborative Filtering--Based Movie
Recommendation System: From-Scratch Implementation, PyTorch
Benchmarking, and Production
Architecture}\label{design-of-a-neural-collaborative-filteringbased-movie-recommendation-system-from-scratch-implementation-pytorch-benchmarking-and-production-architecture}

\textbf{RAHUL K. P}\\
\emph{MSc Computer Science, Final Year}\\
University of Calicut\\
rahul.kp.msc.cs@gmail.com

\textbf{Dr.~Seema S.}\\
\emph{MCA, MBA, MPhil, PhD}\\
University of Calicut\\
seema.karuvarath@gmail.com

\subsection{Abstract}\label{abstract}

The proliferation of digital content platforms has rendered personalized
recommendation systems a foundational component of modern information
retrieval. Classical Matrix Factorization (MF) methods, while
computationally tractable, are fundamentally constrained by the
linearity of the inner product operator, which prevents them from
capturing the non-linear, higher-order dependencies characteristic of
real-world user--item interaction spaces. This paper presents a complete
end-to-end system embodying Neural Collaborative Filtering (NCF),
wherein a Generalized Matrix Factorization (GMF) module and a
Multi-Layer Perceptron (MLP) are fused to model both linear and
non-linear latent factors simultaneously. Two fully isomorphic
implementations are developed: a pedagogical NumPy-based version
featuring hand-derived backpropagation, and an optimized PyTorch version
leveraging Apple Silicon MPS acceleration. Empirical evaluation on the
MovieLens dataset demonstrates that both implementations converge to
equivalent final Binary Cross-Entropy losses (0.2257 and 0.2307,
respectively), while the PyTorch variant achieves a 3.39× speedup in
total training time (3,295 s versus 972 s over 20 epochs). Peak
recommendation quality, measured by Hit Ratio at cutoff 10 (HR@10),
reaches 0.615 for the PyTorch implementation. The system is deployed as
a production-grade microservices architecture comprising a FastAPI
gateway, a dedicated PyTorch inference server, a PostgreSQL persistence
layer, and a Netflix-style frontend with TMDB poster integration. A
hybrid cold-start module supplements the NCF core for new users and
items. The findings validate the feasibility of bridging rigorous
algorithmic pedagogy with industry-standard deployment practices.

\textbf{Index Terms}---Neural Collaborative Filtering, Recommender
Systems, Matrix Factorization, Implicit Feedback, Deep Learning,
Microservices Architecture, Scalable Deployment, MovieLens.

\begin{center}\rule{0.5\linewidth}{0.5pt}\end{center}

\subsection{I. Introduction}\label{i.-introduction}

The exponential expansion of online content catalogues---encompassing
streaming media, e-commerce inventories, and digital knowledge
bases---has rendered the task of surfacing relevant items to individual
users both economically critical and technically challenging.
Recommender systems constitute the principal algorithmic mechanism
through which this filtering is accomplished, and their design has
evolved considerably since the seminal deployment of collaborative
filtering techniques at Usenet and GroupLens in the mid-1990s.

Early approaches to collaborative filtering relied on memory-based
neighbourhood methods that computed user or item similarities directly
from the interaction matrix. While conceptually intuitive, these
techniques suffer from poor scalability to millions of users and items.
Model-based approaches, and most notably Singular Value Decomposition
(SVD)-inspired Matrix Factorization (MF) {[}2{]}, addressed scalability
by projecting users and items into shared low-dimensional latent spaces,
enabling the prediction of unobserved interactions through inner
products of latent vectors. MF remains competitive on explicit feedback
benchmarks; however, when applied to implicit feedback (binary signals
derived from clicks, views, or purchases), the linearity of the inner
product imposes a fundamental representational bottleneck. Specifically,
the inner product is unable to satisfy arbitrary rank orderings of item
relevance, a limitation formally characterised by He et al.~{[}1{]}.

The Neural Collaborative Filtering framework introduced by He et
al.~{[}1{]} resolves this limitation by replacing the inner product with
a neural architecture capable of learning arbitrary continuous functions
of the user and item embeddings. NCF subsumes MF as a special case
through its GMF pathway and augments it with a deep MLP pathway that
captures non-linear feature interactions. Crucially, the two pathways
are fused in a final prediction layer, allowing the model to jointly
leverage the complementary inductive biases of linear factorization and
deep representation learning.

Despite the theoretical importance of the NCF framework, relatively few
published works provide a complete system-level blueprint that connects
the mathematical derivation of NCF to a production-ready deployment.
This paper addresses that gap with three primary contributions. First, a
mathematically rigorous derivation and from-scratch NumPy implementation
of NCF is presented, including explicit manual backpropagation through
the GMF and MLP pathways. Second, an equivalent PyTorch implementation
is benchmarked against the NumPy baseline, providing quantitative
evidence of the computational advantages conferred by automatic
differentiation and hardware-accelerated tensor operations. Third, the
trained NCF model is integrated into a scalable microservices
architecture deployed via Docker, with a Netflix-style frontend and a
hybrid cold-start recommendation module. Taken together, these
contributions form a reproducible blueprint suitable for both academic
research and industrial deployment.

\begin{center}\rule{0.5\linewidth}{0.5pt}\end{center}

\subsection{II. Related Work}\label{ii.-related-work}

Collaborative filtering methods for recommender systems were
systematically formalised by Breese et al.~{[}3{]}, who established the
distinction between memory-based and model-based approaches. The
introduction of probabilistic matrix factorization by Salakhutdinov and
Mnih {[}4{]} and the subsequent deployment of regularised SVD variants
in the Netflix Prize competition {[}2{]} established latent factor
models as the dominant paradigm. Koren et al.~{[}2{]} provided a
comprehensive survey of these methods, demonstrating that MF could be
extended to incorporate temporal dynamics, implicit feedback, and
auxiliary side information.

The application of deep learning to collaborative filtering was
pioneered by several concurrent research threads. Salakhutdinov et
al.~{[}5{]} proposed Restricted Boltzmann Machines for collaborative
filtering, demonstrating that generative probabilistic models could
achieve competitive accuracy. Sedhain et al.~{[}6{]} introduced AutoRec,
which applied autoencoder architectures to reconstruct user or item
rating vectors. Wang et al.~{[}7{]} combined collaborative filtering
with deep content features through a Collaborative Deep Learning
framework. These works collectively established that neural
architectures could be productively applied to the recommendation
problem.

The defining contribution in this area remains the Neural Collaborative
Filtering framework of He et al.~{[}1{]}, which provided both a
theoretical critique of the inner product limitation in MF and a
concrete neural architecture to address it. Subsequent work has extended
NCF in several directions: attention-based variants prioritise
temporally recent interactions {[}8{]}; graph neural network approaches
model higher-order collaborative signals through multi-hop neighbourhood
aggregation {[}9{]}; and self-supervised methods leverage auxiliary
contrastive objectives to improve representation quality under sparse
interaction regimes {[}10{]}.

On the system engineering side, the deployment of recommendation models
at scale has received growing attention. Cheng et al.~{[}11{]} described
the Wide \& Deep architecture deployed at Google Play, which combines a
linear model for memorisation with a deep network for generalisation.
Covington et al.~{[}12{]} detailed the two-stage retrieval and ranking
architecture underpinning YouTube recommendations. These industrial
systems highlight the importance of infrastructure
considerations---including serving latency, cold-start handling, and
feature freshness---that are often underemphasised in academic
treatments. The present work is distinguished by its explicit
integration of these production concerns within the academic NCF
framework, providing a complete system blueprint.

\begin{center}\rule{0.5\linewidth}{0.5pt}\end{center}

\subsection{III. Proposed Methodology}\label{iii.-proposed-methodology}

\subsubsection{A. Mathematical Formulation and Problem
Definition}\label{a.-mathematical-formulation-and-problem-definition}

Let \(U = \{u_1, u_2, \dots, u_m\}\) denote the set of \(m\) users and
\(I = \{i_1, i_2, \dots, i_n\}\) denote the set of \(n\) items. The
user-item interaction matrix \(Y \in \{0, 1\}^{m \times n}\) encodes
implicit feedback, where:

\[y_{ui} = 1 \text{ if interaction } (u,i) \text{ is observed; } y_{ui} = 0 \text{ otherwise.}\]

The recommendation task is formulated as estimating the probability that
user \(u\) will interact with item \(i\), denoted
\(\hat{y}_{ui} \in (0, 1)\). This is equivalent to a binary
classification problem over the user-item pair \((u,i)\). Formally, the
predictive function is:

\[\hat{y}_{ui} = f(u, i \mid \Theta)\]

where \(f\) is a neural interaction function parameterised by
\(\Theta\). The NCF framework realises \(f\) through a fusion of two
complementary sub-networks described below.

\subsubsection{B. Generalized Matrix Factorization
(GMF)}\label{b.-generalized-matrix-factorization-gmf}

The GMF pathway generalises the standard MF inner product by replacing
it with an element-wise (Hadamard) product followed by a learnable
projection. Let \(p_u^G \in \mathbb{R}^K\) and
\(q_i^G \in \mathbb{R}^K\) denote the GMF embedding vectors for user
\(u\) and item \(i\), respectively, where \(K\) is the embedding
dimensionality. The GMF output vector is:

\[\phi^{\text{GMF}} = p_u^G \odot q_i^G\]

where \(\odot\) denotes the Hadamard product. Standard MF is recovered
as a special case when the final projection weight vector is a constant
vector of ones and no non-linear activation is applied. By generalising
the projection to a trainable weight \(h^G \in \mathbb{R}^K\), the GMF
pathway learns a data-adaptive linear combination of the element-wise
interaction features.

\subsubsection{C. Multi-Layer Perceptron (MLP)
Pathway}\label{c.-multi-layer-perceptron-mlp-pathway}

The MLP pathway captures non-linear interaction patterns through a deep
feedforward network. Let \(p_u^M \in \mathbb{R}^K\) and
\(q_i^M \in \mathbb{R}^K\) denote independent MLP embeddings for user
\(u\) and item \(i\). The concatenated input vector to the MLP is:

\[z_1 = [p_u^M ; q_i^M] \in \mathbb{R}^{2K}\]

The MLP applies \(L\) hidden layers with ReLU activations:

\[z_l = \text{ReLU}(W_l^T z_{l-1} + b_l), \quad l = 1, 2, \dots, L\]

where \(W_l \in \mathbb{R}^{d_{l-1} \times d_l}\) and
\(b_l \in \mathbb{R}^{d_l}\) are the weight matrix and bias vector of
layer \(l\), respectively. The output of the final MLP layer,
\(\phi^{\text{MLP}} = z_L\), captures a rich non-linear representation
of the user-item interaction.

\subsubsection{D. NeuMF: Fusion
Architecture}\label{d.-neumf-fusion-architecture}

The complete NeuMF (Neural Matrix Factorization) model concatenates the
outputs of the GMF and MLP pathways and maps the result to a scalar
prediction through a final projection:

\[\hat{y}_{ui} = \sigma(h^T [\phi^{\text{GMF}} ; \phi^{\text{MLP}}])\]

where \(h \in \mathbb{R}^{K + d_L}\) is a trainable weight vector and
\(\sigma(\cdot)\) is the sigmoid activation function, ensuring
\(\hat{y}_{ui} \in (0, 1)\). The separation of GMF and MLP embeddings is
a deliberate design choice: it allows each pathway to learn specialised
representations without constraining the MLP to operate in the same
embedding space as the GMF factorization.

\subsubsection{E. Training Objective and Negative
Sampling}\label{e.-training-objective-and-negative-sampling}

The model is trained by minimising the Binary Cross-Entropy (BCE) loss
over observed positive interactions and sampled negative instances:

\[\mathcal{L} = -\sum_{(u,i) \in Y \cup Y^-} \left[ y_{ui} \cdot \log(\hat{y}_{ui}) + (1 - y_{ui}) \cdot \log(1 - \hat{y}_{ui}) \right]\]

where \(Y\) denotes the set of observed (positive) interactions and
\(Y^-\) is a set of unobserved (negative) samples drawn uniformly at
random from items not interacted with by user \(u\). For each positive
interaction, four negative samples are drawn, yielding a 1:4
positive-to-negative ratio. This negative sampling ratio follows the
empirical recommendation of He et al.~{[}1{]} and was validated in
preliminary experiments on the MovieLens dataset.

\subsubsection{F. Manual Backpropagation in the NumPy
Implementation}\label{f.-manual-backpropagation-in-the-numpy-implementation}

The NumPy implementation derives all gradient computations analytically,
without reliance on automatic differentiation. For the BCE loss with
sigmoid output, the gradient with respect to the pre-activation
prediction score
\(s_{ui} = h^T [\phi^{\text{GMF}} ; \phi^{\text{MLP}}]\) is:

\[\frac{\partial \mathcal{L}}{\partial s_{ui}} = \hat{y}_{ui} - y_{ui}\]

This gradient is then propagated through the concatenation, the MLP
layers via the chain rule, and the element-wise product in the GMF
pathway. For each MLP layer \(l\), the weight gradient is:

\[\frac{\partial \mathcal{L}}{\partial W_l} = z_{l-1} \cdot (\delta_l)^T\]

where
\(\delta_l = \left(\frac{\partial \mathcal{L}}{\partial z_l}\right) \odot \text{ReLU}'(z_l)\)
and \(\delta_{l-1} = W_l \cdot \delta_l\), with \(\text{ReLU}'(x) = 1\)
if \(x > 0\), else \(0\). Embedding gradients are accumulated across all
training samples and applied using Adam optimisation with learning rate
0.001. This explicit derivation constitutes the primary pedagogical
contribution of the NumPy implementation.

\subsubsection{G. Hybrid Cold-Start
Module}\label{g.-hybrid-cold-start-module}

For users and items with insufficient interaction history, a hybrid
recommendation score is computed as a weighted linear combination of the
NCF score and a content/popularity-based cold-start score:

\[\text{Score}(u,i) = \alpha \cdot \text{NCF}(u,i) + (1-\alpha) \cdot \text{ColdStart}(u,i)\]

where \(\alpha \in [0,1]\) is dynamically set based on the number of
interactions available for user \(u\): \(\alpha = 0\) for new users
(cold start mode), transitioning linearly to \(\alpha = 1\) after a
configurable interaction threshold (default: 10 interactions). The
cold-start score combines item popularity (global interaction frequency)
with content-based genre similarity, enabling meaningful recommendations
from the first user session.

\begin{center}\rule{0.5\linewidth}{0.5pt}\end{center}

\subsection{IV. System Design and
Architecture}\label{iv.-system-design-and-architecture}

\subsubsection{A. Microservices
Decomposition}\label{a.-microservices-decomposition}

The production architecture decomposes the recommendation system into
independent microservices, each with a single responsibility,
communicating over HTTP/REST. This decomposition ensures that
compute-intensive model inference does not contend with the main
application thread, enabling horizontal scaling of individual services
in response to differential load patterns. The four primary services
are: (1) the API Gateway, (2) the Model Inference Server, (3) the Data
Persistence Service, and (4) the Frontend Service.

\subsubsection{B. API Gateway (FastAPI)}\label{b.-api-gateway-fastapi}

The API Gateway is implemented using FastAPI, a modern Python ASGI
(Asynchronous Server Gateway Interface) framework that supports fully
asynchronous request handling via Python's asyncio runtime. The gateway
exposes four primary endpoint groups:

\begin{itemize}
\tightlist
\item
  \texttt{/recommend} for personalised recommendation retrieval;
\item
  \texttt{/search} for query-driven item lookup with NCF refinement;
\item
  \texttt{/interact} for recording user interaction events that trigger
  model score updates; and
\item
  \texttt{/compare} for serving the benchmark dashboard data.
\end{itemize}

Response latency for recommendation queries is maintained below 100 ms
through in-process Redis caching of frequently requested user
recommendation lists, with a cache eviction time-to-live of 300 seconds.

\subsubsection{C. Model Inference
Server}\label{c.-model-inference-server}

The PyTorch NCF model is hosted in an isolated inference service that
loads the trained model checkpoint at startup and exposes a
\texttt{/predict} endpoint accepting batches of
\texttt{(user\_id,\ item\_id)} pairs. Isolation of the inference service
prevents GPU/MPS memory contention with the API Gateway process and
allows the inference service to be independently scaled or replaced
without modifying the gateway. Inference batching is supported to
amortise per-request overhead: recommendation lists of top-\(K\) items
are generated by computing predictions over all candidate items in a
single forward pass and applying argmax-\(K\) selection.

\subsubsection{D. Data Persistence
Layer}\label{d.-data-persistence-layer}

A PostgreSQL relational database stores three primary entity types:
users (\texttt{user\_id}, registration timestamp, demographic metadata),
movies (\texttt{movie\_id}, title, genre vector, TMDB metadata), and
interactions (\texttt{user\_id}, \texttt{movie\_id},
\texttt{interaction\_type} \(\in \{\text{click}, \text{like}\}\),
timestamp, \texttt{session\_id}). Genre metadata is stored as a
PostgreSQL \texttt{ARRAY} column and indexed for efficient content-based
filtering. Interaction records are append-only to preserve the temporal
sequence required for online learning. Redis is deployed as a secondary
cache layer, storing serialised recommendation lists keyed by
\texttt{user\_id}, and invalidating entries upon receipt of a new
interaction event from the corresponding user.

\subsubsection{E. Containerisation and
Orchestration}\label{e.-containerisation-and-orchestration}

The complete system is containerised using Docker, with a Docker Compose
specification defining five named services: \texttt{api-gateway},
\texttt{model-server}, \texttt{postgres}, \texttt{redis}, and
\texttt{frontend}. Multi-stage Docker builds are employed to minimise
production image sizes: the \texttt{model-server} image, for instance,
separates a build stage (installing PyTorch and dependencies) from a
runtime stage (copying only the required Python environment and model
artefacts), reducing the final image size by approximately 60\% relative
to a single-stage build. Environment-specific configuration is managed
via \texttt{.env} files, ensuring strict separation between development
and production credentials.

\subsubsection{F. Netflix-Style
Frontend}\label{f.-netflix-style-frontend}

The frontend presents a Netflix-inspired poster grid interface with
three primary sections: Home (populated by cold-start recommendations
for unauthenticated or new users), Trending (populated by global
popularity scores), and Recommended (populated by NCF personalised
scores for returning users). Movie posters are retrieved from the TMDB
API using each film's TMDB identifier. User interaction events (clicks
and likes) are captured by the frontend and transmitted asynchronously
to the \texttt{/interact} endpoint, triggering NCF score recomputation
for the affected user in the background. The search interface invokes
the \texttt{/search} endpoint; when the queried title is found in the
database, NCF-refined results are returned alongside the direct match,
surfacing contextually related films.

\begin{center}\rule{0.5\linewidth}{0.5pt}\end{center}

\subsection{V. Experimental Setup}\label{v.-experimental-setup}

\subsubsection{A. Dataset}\label{a.-dataset}

Experiments were conducted on the MovieLens 1M dataset {[}13{]}, which
contains 1,000,209 ratings from 6,040 users across 3,706 movies.
Following standard practice for implicit feedback modelling {[}1{]},
explicit ratings were binarised: any rating of 1 or above was treated as
a positive interaction (\(y_{ui} = 1\)). To ensure evaluation quality,
only users with at least 20 interactions were retained. The
leave-one-out evaluation protocol was applied: for each user, the most
recent interaction was reserved for testing, the penultimate interaction
for validation, and all remaining interactions for training.

\subsubsection{B. Negative Sampling}\label{b.-negative-sampling}

For each positive test interaction, 99 negative items were sampled
uniformly at random from the set of items not interacted with by the
test user, consistent with the evaluation protocol of He et al.~{[}1{]}.
Model performance is thus assessed over a ranked list of 100 items (one
positive, 99 negatives), enabling computation of ranking-based metrics.
During training, four negative samples were generated per positive
interaction per epoch, yielding a training set of approximately 5
million instances per epoch.

\subsubsection{C. Model Configuration and
Hyperparameters}\label{c.-model-configuration-and-hyperparameters}

Both the NumPy and PyTorch implementations share an identical
architecture: embedding dimensionality \(K = 32\) for both GMF and MLP
pathways; MLP hidden layer dimensions \([256, 128, 64]\) with ReLU
activations; batch size 256; learning rate 0.001 with Adam optimisation;
and 20 training epochs. The final output dimension of the MLP is 64, so
the concatenated fusion vector
\([\phi^{\text{GMF}} ; \phi^{\text{MLP}}]\) has dimension 96. Weight
initialisation follows a normal distribution with mean 0 and standard
deviation 0.01 for embedding matrices, and Xavier uniform initialisation
for MLP weight matrices.

\subsubsection{D. Evaluation Metrics}\label{d.-evaluation-metrics}

Performance is quantified using two standard ranking metrics. Hit Ratio
at \(K\) (HR@\(K\)) measures whether the positive test item appears in
the top-\(K\) ranked recommendations for the user, averaged across all
test users. For \(K = 10\):

\[\text{HR@10} = \frac{|\{u : \text{rank}(i_u^+) \le 10\}|}{|U|}\]

where \(\text{rank}(i_u^+)\) is the rank of the positive item for user
\(u\). Normalised Discounted Cumulative Gain at \(K\) (NDCG@\(K\))
additionally penalises recommendations that rank the positive item lower
within the top-\(K\) list, providing a position-sensitive quality
measure. Both metrics range from 0 to 1, with higher values indicating
better recommendation quality.

\subsubsection{E. Hardware and Software
Environment}\label{e.-hardware-and-software-environment}

All experiments were conducted on an Apple M1 Pro system with 32 GB
unified memory. The PyTorch implementation utilised the Metal
Performance Shaders (MPS) backend for GPU-accelerated tensor operations
on the Apple Silicon neural engine. The NumPy implementation executed on
CPU only, with no vectorisation optimisations beyond NumPy's default
BLAS routines. Software versions: Python 3.11.4, PyTorch 2.1.0, NumPy
1.25.2, FastAPI 0.104.0, PostgreSQL 15.2.

\begin{center}\rule{0.5\linewidth}{0.5pt}\end{center}

\subsection{VI. Results and Performance
Evaluation}\label{vi.-results-and-performance-evaluation}

\subsubsection{A. Training Loss
Convergence}\label{a.-training-loss-convergence}

Figure 1 presents the BCE training loss trajectories for both
implementations over 20 epochs. Both curves originate from the same
initialisation point (\(\text{BCE} \approx 0.365\) at epoch 1) and
converge monotonically, confirming that the manual backpropagation
implementation correctly computes gradients and that both
implementations are solving the same optimisation problem. The Scratch
NCF achieves a final loss of 0.2257, marginally lower than the PyTorch
NCF final loss of 0.2307---a difference of 0.005 that is attributable to
minor numerical precision differences between Python float64 (NumPy) and
float32 (PyTorch) arithmetic. Both losses substantially undercut the
random baseline of 0.693, confirming that both implementations have
learned meaningful user--item representations.

\end{document}
